\section{Experimental Setup}
\label{sec:setup}

\begin{table}[t]
\centering\small
\begin{adjustbox}{max width=\linewidth}
\begin{tabular}{l l r l r l}
\toprule
Texts & Written by & Prompts & Genres & Texts & Formats judged \\
\midrule
Chat-app essays & 5 frontier models, chat apps & 38 & essays & 190 & lineup, one text, yes/no \\
API texts & 5 frontier models, API & 120 & all four & 600 & lineup, one text \\
Open-weight texts & 4 open-weight models & 120 & all four & 480 & lineup, one text, yes/no \\
\bottomrule
\end{tabular}
\end{adjustbox}
\caption{The three sets of texts. The 38 chat-app prompts are among the 120;
all 120 are in the repository (\texttt{data/prompts\_v2.csv}), with one
example per genre in Appendix~\ref{app:data}. In each set, the models that wrote the texts also judge them; the chat-app essays carry the product labels listed in Section~\ref{sec:setup}.}
\label{tab:data}
\end{table}


\subsection{Texts}

We use three sets of texts, summarised in Table~\ref{tab:data}. All are built
from one bank of prompts designed so that answers to the same prompt differ
mainly in style: each prompt fixes the topic and the points the answer must
cover, and essay prompts ask for plain paragraphs with no lists, headers or
bold (Appendix~\ref{app:data} gives an example of each genre).

\paragraph{Chat-app essays} Five frontier models, GPT-5.5, Claude Opus
4.7, Gemini 3.5 Flash, Grok 4.3 and DeepSeek's chat model, each answered 39
argumentative prompts in their consumer chat apps at default settings,
between 15 June and 31 July 2026. One prompt was dropped because two of its
answers were corrupted copies, leaving 38 prompts and 190 essays.

\paragraph{API texts} To test whether results hold beyond one small set of
essays, the same five models answered 120 prompts through the API on 1
October 2026 (temperature 0.7, no system prompt): 60 essays (including the 38
above), 20 Python functions, 20 short explanations and 20 math word problems,
giving 600 texts.

\paragraph{Open-weight texts} Four open-weight models, Qwen2.5-7B
(Alibaba), Llama-3.1-8B (Meta), Mistral-7B-v0.3 (Mistral AI) and Phi-4-mini
(Microsoft), answered the same 120 prompts (temperature 0.7), giving 480
texts, plus a second, independent answer to every prompt. We generated these
texts ourselves from public weights, so none carries a watermark.

\subsection{Judges}

In each set, the models that wrote the texts are the judges, so ``self''
always means the same model on both sides. The frontier models judge through
OpenRouter at temperature~0.\footnote{Six of Claude's 38 chat-app essays were
written by Claude Sonnet 4.6 rather than Opus 4.7; dropping them leaves
Claude's lineup self-naming rate at 27 of 32 (84\%). The DeepSeek endpoint
could not be pinned to a snapshot.} The open-weight models run in 4-bit on a
single laptop GPU with greedy decoding (Appendix~\ref{app:openweight}). Judges
are never given examples, and see the prompt only in the answer-first
condition below. With five candidate
authors chance is 20\%; with four, 25\%.

\subsection{Question Formats}

Judges answer in the following formats; the exact instructions are in
Appendix~\ref{app:instructions}.
\begin{itemize}
\item \textbf{Lineup.} The judge sees all answers to one prompt, labelled A,
B, C, \ldots, and names an author for each. It is told it may name the same
model more than once, and the order of answers is balanced across judges.
\item \textbf{One text at a time.} The judge sees a single answer and names
its author from the panel.
\item \textbf{Yes/no.} The judge is asked ``Did you write this text?'' in a
fresh conversation. For open-weight judges we read the probability of
``yes''; no frontier provider returned probabilities, so for frontier judges
we use the answer itself.
\item \textbf{Answer first (open-weight only).} Following
\citet{khullar2026selfattribution}, the judge first answers the prompt
itself and is then asked the yes/no question about either that same answer, a
separate answer of its own, or a peer's answer.
\item \textbf{Shuffled lineup (frontier, chat-app essays only).} The lineup
again, with every essay's words shuffled, which keeps the vocabulary but
removes word order and structure.
\end{itemize}

\subsection{A Reference Classifier}

To measure how much the texts reveal about their authors, we train a random
forest on 18 simple style features (length, punctuation, vocabulary richness,
readability), always testing on prompts it was not trained on. It identifies
the author of \EngineAAcc{} of the chat-app essays (95\% CI \EngineACI{};
chance \Chance{}) and \OWClfPooled{} of the open-weight texts (chance 25\%),
more than any judge achieves (Appendix~\ref{app:classifier}).
